\documentclass{report}
\usepackage{luatexja} % LuaTeXで日本語を使うためのパッケージ
\usepackage{luatexja-fontspec} % LuaTeX用の日本語フォント設定

% --- 数学関連 ---
\usepackage{amsmath, amssymb, amsfonts, mathtools, bm, amsthm} % 基本的な数学パッケージ
\usepackage{type1cm, upgreek} % 数式フォントとギリシャ文字
\usepackage{physics, mhchem} % 物理や化学の記号や式の表記を簡単にする

% --- 表関連 ---
\usepackage{multirow, longtable, tabularx, array, colortbl, dcolumn, diagbox} % 表のレイアウトを柔軟にする
\usepackage{tablefootnote, truthtable} % 表中に注釈を追加、真理値表
\usepackage{tabularray} % 高度な表組みレイアウト

% --- グラフィック関連 ---
\usepackage{tikz, graphicx} % 図の描画と画像の挿入
\usepackage{background} % ウォーターマークの設定
\usepackage{caption, subcaption} % 図や表のキャプション設定
\usepackage{float, here} % 図や表の位置指定

% --- レイアウトとページ設定 ---
\usepackage{fancyhdr} % ページヘッダー、フッター、余白の設定
\usepackage[top = 20truemm, bottom = 20truemm, left = 20truemm, right = 20truemm]{geometry}
\usepackage{fancybox, ascmac} % ボックスのデザイン

% --- 色とスタイル ---
\usepackage{xcolor, color, colortbl, tcolorbox} % 色とカラーボックス
\usepackage{listings, jvlisting} % コードの色付けとフォーマット

% --- 参考文献関連 ---
\usepackage{biblatex, usebib} % 参考文献の管理と挿入
\usepackage{url, hyperref} % URLとリンクの設定

% --- その他の便利なパッケージ ---
\usepackage{footmisc} % 脚注のカスタマイズ
\usepackage{multicol} % 複数段組
\usepackage{comment} % コメントアウトの拡張
\usepackage{siunitx} % 単位の表記
\usepackage{docmute}
% \usepackage{appendix}
% --- tcolorboxとtikzの設定 ---
\tcbuselibrary{theorems, breakable} % 定理のボックスと改ページ設定
\usetikzlibrary{decorations.markings, arrows.meta, calc} % tikzの装飾や矢印の設定

% --- 定理スタイルと数式設定 ---
\theoremstyle{definition} % 定義スタイル
\numberwithin{equation}{section} % 式番号をサブセクション単位でリセット

% --- hyperrefの設定 ---
\hypersetup{
  setpagesize = false,
  bookmarks = true,
  bookmarksdepth = tocdepth,
  bookmarksnumbered = true,
  colorlinks = false,
  pdftitle = {}, % PDFタイトル
  pdfsubject = {}, % PDFサブジェクト
  pdfauthor = {}, % PDF作者
  pdfkeywords = {} % PDFキーワード
}

% --- siunitxの設定 ---
\sisetup{
  table-format = 1.5, % 小数点以下の桁数
  table-number-alignment = center, % 数値の中央揃え
}

% --- 透かし画像の設定 --- 
\backgroundsetup{
  scale=0.5,                       % 画像のスケール
  % color=black,                   % 画像の色（透かし用に半透明が推奨）
  opacity=0.2,                   % 透かしの透明度（0が完全透明、1が完全不透明）
  angle=0,                       % 画像の角度
  position = current page.south east,  % ページの右下
  hshift=-6cm, % 右方向へのシフト（負の値で内側に移動）
  vshift=5cm,  % 上方向へのシフト（正の値で内側に移動）
  contents={\includegraphics{./fig/appilogo-circular-full.png}} % 画像のパス
}

% --- その他の設定 ---
\allowdisplaybreaks % 数式の途中改ページ許可
\newcolumntype{t}{!{\vrule width 0.1pt}} % 新しいカラムタイプ
\newcolumntype{b}{!{\vrule width 1.5pt}} % 太いカラム
\UseTblrLibrary{amsmath, booktabs, counter, diagbox, functional, hook, html, nameref, siunitx, varwidth, zref} % tabularrayのライブラリ
\setlength{\columnseprule}{0.4pt} % カラム区切り線の太さ
\captionsetup[figure]{font = bf} % 図のキャプションの太字設定
\captionsetup[table]{font = bf} % 表のキャプションの太字設定
\captionsetup[lstlisting]{font = bf} % コードのキャプションの太字設定
\captionsetup[subfigure]{font = bf, labelformat = simple} % サブ図のキャプション設定
\setcounter{secnumdepth}{4} % セクションの深さ設定
\newcolumntype{d}{D{.}{.}{5}} % 数値のカラム
\newcolumntype{M}[1]{>{\centering\arraybackslash}m{#1}} % センター揃えのカラム
\DeclareMathOperator{\diag}{diag}
\everymath{\displaystyle} % 数式のスタイル
\newcommand{\inner}[2]{\left\langle #1, #2 \right\rangle}
\renewcommand{\figurename}{図}
\renewcommand{\i}{\mathrm{i}} % 複素数単位i
\renewcommand{\laplacian}{\grad^2} % ラプラシアンの記号
\renewcommand{\thesubfigure}{(\alph{subfigure})} % サブ図の番号形式
\newcommand{\m}[3]{\multicolumn{#1}{#2}{#3}} % マルチカラムのショートカット
\renewcommand{\r}[1]{\mathrm{#1}} % mathrmのショートカット
\newcommand{\e}{\mathrm{e}} % 自然対数の底e
\newcommand{\Ef}{E_{\mathrm{F}}} % フェルミエネルギー
\renewcommand{\c}{\si{\degreeCelsius}} % 摂氏記号
\renewcommand{\d}{\r{d}} % d記号
\renewcommand{\t}[1]{\texttt{#1}} % タイプライタフォント
\newcommand{\kb}{k_{\mathrm{B}}} % ボルツマン定数
% \renewcommand{\phi}{\varphi} % ϕをφに変更
\renewcommand{\epsilon}{\varepsilon}
\newcommand{\fullref}[1]{\textbf{\ref{#1} \nameref{#1}}}
\newcommand{\reff}[1]{\textbf{図\ref{#1}}} % 図参照のショートカット
\newcommand{\reft}[1]{\textbf{表\ref{#1}}} % 表参照のショートカット
\newcommand{\refe}[1]{\textbf{式\eqref{#1}}} % 式参照のショートカット
\newcommand{\refp}[1]{\textbf{コード\ref{#1}}} % コード参照のショートカット
\renewcommand{\lstlistingname}{コード} % コードリストの名前
\renewcommand{\theequation}{\thesection.\arabic{equation}} % 式番号の形式
\renewcommand{\footrulewidth}{0.4pt} % フッターの線
\newcommand{\mar}[1]{\textcircled{\scriptsize #1}} % 丸囲み文字
\newcommand{\combination}[2]{{}_{#1} \mathrm{C}_{#2}} % 組み合わせ
\newcommand{\thline}{\noalign{\hrule height 0.1pt}} % 細い横線
\newcommand{\bhline}{\noalign{\hrule height 1.5pt}} % 太い横線

% --- カスタム色定義 ---
\definecolor{burgundy}{rgb}{0.5, 0.0, 0.13} % バーガンディ色
\definecolor{charcoal}{rgb}{0.21, 0.27, 0.31} % チャコール色
\definecolor{forest}{rgb}{0.0, 0.35, 0} % 森の緑色

% --- カスタム定理環境の定義 ---
\newtcbtheorem[number within = chapter]{myexc}{練習問題}{
  fonttitle = \gtfamily\sffamily\bfseries\upshape,
  colframe = forest,
  colback = forest!2!white,
  rightrule = 1pt,
  leftrule = 1pt,
  bottomrule = 2pt,
  colbacktitle = forest,
  theorem style = standard,
  breakable,
  arc = 0pt,
}{exc-ref}
\newtcbtheorem[number within = chapter]{myprop}{命題}{
  fonttitle = \gtfamily\sffamily\bfseries\upshape,
  colframe = blue!50!black,
  colback = blue!50!black!2!white,
  rightrule = 1pt,
  leftrule = 1pt,
  bottomrule = 2pt,
  colbacktitle = blue!50!black,
  theorem style = standard,
  breakable,
  arc = 0pt
}{proposition-ref}
\newtcbtheorem[number within = chapter]{myrem}{注意}{
  fonttitle = \gtfamily\sffamily\bfseries\upshape,
  colframe = yellow!20!black,
  colback = yellow!50,
  rightrule = 1pt,
  leftrule = 1pt,
  bottomrule = 2pt,
  colbacktitle = yellow!20!black,
  theorem style = standard,
  breakable,
  arc = 0pt
}{remark-ref}
\newtcbtheorem[number within = chapter]{myex}{例題}{
  fonttitle = \gtfamily\sffamily\bfseries\upshape,
  colframe = black,
  colback = white,
  rightrule = 1pt,
  leftrule = 1pt,
  bottomrule = 2pt,
  colbacktitle = black,
  theorem style = standard,
  breakable,
  arc = 0pt
}{example-ref}
\newtcbtheorem[number within = chapter]{exc}{Requirement}{myexc}{exc-ref}
\newcommand{\rqref}[1]{{\bfseries\sffamily 練習問題 \ref{exc-ref:#1}}}
\newtcbtheorem[number within = chapter]{definition}{Definition}{mydef}{definition-ref}
\newcommand{\dfref}[1]{{\bfseries\sffamily 定義 \ref{definition-ref:#1}}}
\newtcbtheorem[number within = chapter]{prop}{命題}{myprop}{proposition-ref}
\newcommand{\prref}[1]{{\bfseries\sffamily 命題 \ref{proposition-ref:#1}}}
\newtcbtheorem[number within = chapter]{rem}{注意}{myrem}{remark-ref}
\newcommand{\rmref}[1]{{\bfseries\sffamily 注意 \ref{remark-ref:#1}}}
\newtcbtheorem[number within = chapter]{ex}{例題}{myex}{example-ref}
\newcommand{\exref}[1]{{\bfseries\sffamily 例題 \ref{example-ref:#1}}}
% --- 再定義コマンド ---
% \mathtoolsset{showonlyrefs=true} % 必要な式番号のみ表示
\pagestyle{fancy} % ヘッダー・フッターのスタイル設定
\chead{応用量子物性講義ノート} % 中央ヘッダー
% \rhead{}
\fancyhead[R]{\rightmark}
\renewcommand{\sectionmark}[1]{\markright{\thesection\ #1}}
\cfoot{\thepage} % 中央フッターにページ番号
\lhead{}
\rfoot{Yuto Masuda and Haruki Aoki} % 右フッターに名前
\setcounter{tocdepth}{4} % 目次の深さ
\makeatletter
\@addtoreset{equation}{section} % サブセクションごとに式番号をリセット
\makeatother

% --- メタ情報 ---
\title{応用量子物性講義ノート}
\date{更新日\today}
\author{Yuto Masuda and Haruki Aoki}

\begin{document}
  \begin{itembox}[l]{変分法の基本原理}
  任意の状態ベクトル$\ket{\psi}$に対して$\ket{\psi}$でのエネルギー関数$E(\psi)$について，
    \begin{align*}
      E(\psi) = \cfrac{\mel**{\psi}{\hat{H}}{\psi}}{\braket{\psi}} \geq E_0 \label{minimum}
    \end{align*}
  なる不等式が成り立つ．
  ただし$E_0$は$\hat{H}$の固有エネルギーの中で最低のものである．
  \end{itembox}
  \begin{itembox}[l]{1次摂動によるエネルギー補正}
    \begin{align*}
      E_n^{(1)} = \ev**{\hat{V}}{n^{(0)}}
    \end{align*}
  \end{itembox}
  \begin{itembox}[l]{1次摂動による固有ベクトル補正}
    \begin{align*}
      \ket{n^{(1)}} = \sum_{m\neq n}\cfrac{\mel**{m^{(0)}}{\hat{V}}{n^{(0)}}}{E_n^{(0)} - E_m^{(0)}}\ket{m^{(0)}}
    \end{align*}
  \end{itembox}
  \begin{itembox}[l]{2次摂動によるエネルギー補正}
    \begin{align*}
      E_n^{(2)} = \sum_{m\neq n}\cfrac{\abs{\mel**{m^{(0)}}{\hat{V}}{n^{(0)}}}^2}{E_n^{(0)} - E_m^{(0)}}
    \end{align*}
  \end{itembox}
  \begin{itembox}[l]{縮退がある場合の摂動論による1次エネルギー補正}
    \begin{align*}
      E_n^{(1)} = \frac{1}{2}\qty[(V_{aa} + V_{bb}\pm\sqrt{(V_{aa} - V_{bb})^2 + 4\abs{V_{ab}}^2})]
    \end{align*}
  \end{itembox}
  \begin{itembox}[l]{相互作用表示}
    \begin{align*}
      \ket{\psi(t)}_{\r{I}} \coloneqq \exp\qty(\i\frac{\hat{H}^{(0)}}{\hbar}t)\ket{\psi(t)}
    \end{align*}
  \end{itembox}
  \begin{itembox}[l]{朝永・Schwinger方程式}
    \begin{align*}
      \begin{dcases}
        \i\hbar\dv{t}\ket{\psi(t)}_{\r{I}} = \hat{V}_{\r{I}}(t)\ket{\psi(t)}_{\r{I}} \\
        \hat{V}_{\r{I}}(t) = \exp\qty(\i\frac{\hat{H}^{(0)}}{\hbar}t)\hat{V}(t)\exp\qty(-\i\frac{\hat{H}^{(0)}}{\hbar}t)
      \end{dcases}
    \end{align*}
  \end{itembox}
  \begin{itembox}[l]{非定常摂動量子系の時間発展}
    \begin{align*}
      \i\hbar\dv{t}c_m(t) &= \sum_{n}c_n(t)V_{mn}(t)\e^{\i\omega_{mn}t}
    \end{align*}
  \end{itembox}
  \begin{itembox}[l]{共鳴条件}
    \begin{align*}
      \omega = \omega_{21} = \frac{E_2 - E_1}{\hbar}
    \end{align*}
  \end{itembox}
  \begin{itembox}[l]{$\ket{\psi(t_0)}_{\r{I}} = \ket{i}$の時間発展}
    \begin{align*}
      \begin{dcases}
        \ket{\psi(t)}_{\r{I}}& = \qty(1 + c_{i, i}^{(1)}(t))\ket{i} + \sum_{n\neq i}c_{n, i}^{(1)}(t)\ket{n} \\
        c_{n, i}^{(1)}(t) &= -\frac{\i}{\hbar}\int_{t_0}^{t}V_{n, i}\e^{i\omega_{ni}t}\dd{t}
      \end{dcases}
    \end{align*}
  \end{itembox}
  \begin{itembox}[l]{Fermiの黄金律}
    \begin{align*}
      \omega_{i\to f} = \frac{2\pi}{\hbar}\abs{\mel**{f}{\hat{V}}{i}}^2\delta(E_f - E_i)
    \end{align*}
  \end{itembox}
  \begin{itembox}[l]{電磁場中の電子のハミルトニアン}
    \begin{align*}
      H = \frac{1}{2m}\qty(\bm{p} + e\bm{A})^2 - e\phi
    \end{align*}
  \end{itembox}
  \begin{itembox}[l]{局所ゲージ変換に対して不変なSchrödinger方程式}
    \begin{align*}
      \i\hbar\pdv{t}\psi = \qty[\frac{1}{2m}\qty(\hat{\bm{p}} + e\hat{\bm{A}})^2 - e\phi]\psi
    \end{align*}
  \end{itembox}
  \begin{itembox}[l]{散乱問題の境界条件}
    \begin{align*}
      \psi(\bm{r}) = \e^{\i kz} + f(\theta) \frac{\e^{\i kr}}{r}
    \end{align*}
  \end{itembox}
  \begin{itembox}[l]{散乱振幅と微分断面積の関係}
    \begin{align*}
      \sigma(\theta) = \abs{f(\theta)}^2
    \end{align*}
  \end{itembox}
  \begin{itembox}[l]{球対称ポテンシャルの散乱振幅}
    \begin{align*}
      f^{(1)}(\theta) =  -\frac{2m}{\hbar^2 q} \int_{0}^{\infty} rV(r)\sin qr \dd{r}
    \end{align*}
  \end{itembox}
  \begin{itembox}[l]{部分波展開した散乱の波動関数}
    \begin{align*}
      \psi(\bm{r}) = \sum_{l=0}^{\infty} (2l+1)\i^l j_l(kr)P_l(\cos\theta) + \sum_{l=0}^{\infty} (2l+1)a_l P_l(\cos\theta) \frac{\e^{\i kr}}{r}
    \end{align*}
  \end{itembox}
  \begin{itembox}[l]{光学定理}
    \begin{align*}
      \sigma^{\r{tot}} = \frac{4\pi}{k} \Im f(0)
    \end{align*}
  \end{itembox}
  \begin{itembox}[l]{特殊相対性原理}
    あらゆる慣性系で同じ物理法則が成り立つ．
  \end{itembox}
  \begin{itembox}[l]{光速度不変の原理}
    あらゆる慣性形で真空中の光の速さは同一である．
  \end{itembox}
  \begin{itembox}[l]{Lorentz変換}
    \begin{align*}
      \begin{pmatrix}
        ct' \\ x'
      \end{pmatrix}
      =
      \frac{1}{\sqrt{1 - (v/c)^2}}
      \begin{pmatrix}
        1 & -v/c\\
        -v/c & 1
      \end{pmatrix}
      \begin{pmatrix}
        ct\\x
      \end{pmatrix}
    \end{align*}
  \end{itembox}
  \begin{itembox}[l]{速度の合成}
    \begin{align*}
      V = \frac{v + u'}{1 + \frac{vu'}{c^2}}
    \end{align*}
  \end{itembox}
  \begin{itembox}[l]{Lorentz収縮}
    \begin{align*}
      L' = \sqrt{1-\qty(\frac{v}{c})^2}L
    \end{align*}
  \end{itembox}
  \begin{itembox}[l]{時間の遅れ}
    \begin{align*}
      \Delta T = \frac{1}{\sqrt{1 - (v/c)^2}} \Delta \tau
    \end{align*}
  \end{itembox}
  \begin{itembox}[l]{Lorentz変換に対して共変なMaxwell方程式}
    \begin{align*}
      &\partial_\mu F_{\nu\lambda} + \partial_\nu F_{\lambda\mu} + \partial_\lambda F_{\mu\nu} = 0\\
      &\partial^{\nu}H_{\nu\mu}=j_\mu\\
      &\partial_\mu j^\mu = 0
    \end{align*}
  \end{itembox}
  \begin{itembox}[l]{Klein-Gordon方程式}
    \begin{align*}
      \qty(\frac{1}{c^2}\frac{\partial^2}{\partial t^2} - \grad^2 + \qty(\frac{mc}{\hbar})^2)\psi = 0
    \end{align*}
  \end{itembox}
  \begin{itembox}[l]{Dirac方程式}
    \begin{align*}
      \qty(\i\hbar \frac{\partial}{\partial t} + \i \hbar c \bm{\alpha}\cdot \grad - \beta mc^2)\psi = 0
    \end{align*}
  \end{itembox}
\end{document}